\documentclass[11pt,a4paper]{article}
\usepackage[margin=2.5cm]{geometry}
\usepackage{hyperref}
\usepackage{amsmath,amssymb}
\usepackage{graphicx}
\usepackage{enumitem}
\usepackage{booktabs}
\usepackage{xcolor}
\usepackage{listings}
\usepackage{tcolorbox}
\usepackage{titlesec}

\hypersetup{colorlinks=true, linkcolor=blue, urlcolor=blue, citecolor=blue}

\lstset{
  language=Python,
  basicstyle=\ttfamily\small,
  keywordstyle=\color{blue}\bfseries,
  commentstyle=\color{gray},
  stringstyle=\color{red},
  breaklines=true,
  frame=single,
  numbers=left,
  numberstyle=\tiny\color{gray},
}

\title{%
  \textbf{Meeting Memo \#2} \\[6pt]
  \large Objective 1: Physics-Informed XAI Transformer \\
  for the Soil--Foundation--OWT $f_0$ Correlation Chain
}
\author{Youssef --- PFE Research Project}
\date{February 20, 2026}

\begin{document}
\maketitle

% ============================================================
\section*{Pre-Meeting Note: Remote Desktop Setup for John}
% ============================================================

\begin{tcolorbox}[colback=blue!5, colframe=blue!50!black, title=IT Support Note]
Before the research discussion, I helped \textbf{John} connect his personal computer
to the \textbf{university desktop} via the \textbf{Remote Desktop} Windows app.

\textbf{Problem:} The two computers were on different Wi-Fi networks, meaning they
had different IP address ranges (e.g., John's laptop was on \texttt{192.168.1.x}
at home while the university PC was on \texttt{10.0.5.x}). Windows Remote Desktop
requires both machines to be reachable on the same network --- when IP addresses
start with different numbers, direct connection is impossible.

\textbf{Solution:} We installed \textbf{Tailscale} (\url{https://tailscale.com}) on
both computers. Tailscale is a zero-config VPN built on WireGuard that assigns
every device a shared virtual IP address starting with the same prefix
(\texttt{100.x.x.x}). This means both machines appear to be on the same local network,
regardless of their actual Wi-Fi. Steps taken:

\begin{enumerate}[noitemsep]
  \item Installed Tailscale on John's laptop and the university desktop.
  \item Logged in with the same Tailscale account (or shared via Tailscale admin).
  \item Both machines received IPs like \texttt{100.64.x.x} --- same prefix.
  \item In the Remote Desktop app, John typed the university PC's \textbf{Tailscale IP}
        (e.g., \texttt{100.64.0.2}) instead of the old LAN IP.
  \item Connection succeeded instantly --- secure, encrypted, no port forwarding needed.
\end{enumerate}

\textbf{Key takeaway:} For Remote Desktop to work between two PCs, their IPs must be
routable to each other (i.e., appear to start with the same network prefix).
Tailscale creates this shared virtual network transparently.
\end{tcolorbox}

\vspace{1em}
\hrule
\vspace{1em}

% ============================================================
\section{Context: What Is Objective 1?}
% ============================================================

From the project proposal (Guangdong Provincial Natural Science Fund):

\begin{quote}
\textit{``Characterize how typhoon-induced soil stiffness degradation in Shenzhen
marine clays alters monopile stiffness, and how these changes drive OWT system
natural frequency shifts ($f_0$). Train a physics-informed XAI transformer that
unifies soil, foundation, and superstructure response modelling. Utilize attention
roll-out to quantify and clarify the multi-stage correlation chain.''}
\end{quote}

\textbf{Sources:}
\begin{itemize}[noitemsep]
  \item Abstract Summary Plan ENG (7).pdf, Objective~(1).
  \item PostTyphoon\_OWT\_natural\_frequency\_shift\_prediction.pdf, Section~II-1, p.\,18.
\end{itemize}

The \textbf{correlation chain} to capture is:

\[
\underbrace{G/G_{\max} \downarrow}_{\text{Soil stiffness degradation}}
\;\longrightarrow\;
\underbrace{K_R, K_L, K_{LR} \downarrow}_{\text{Foundation stiffness change}}
\;\longrightarrow\;
\underbrace{f_0^{\text{OWT}} \downarrow}_{\text{Natural frequency shift}}
\]

% ============================================================
\section{Step-by-Step Implementation Plan}
% ============================================================

% ------------------------------------------------------------
\subsection{Step 1: Understand the Physics (the TIM Formulation)}
% ------------------------------------------------------------

The Thermodynamically consistent Inertial Macroelement (TIM) by Gorini (2025)
provides the mathematical blueprint. Foundation stiffness degrades as:

\begin{equation}
  H_{ij}(t) = H^0_{ij} - \sum_{n=1}^{20} w_n(t) \cdot H^n_{ij}
  \label{eq:stiffness_degradation}
\end{equation}

where:
\begin{itemize}[noitemsep]
  \item $H^0_{ij}$: initial 3DOF stiffness matrix (lateral, rotational, coupled),
  \item $H^n_{ij}$: plastic stiffness drop matrix for the $n$-th yield surface,
  \item $w_n(t) \in [0,1]$: sigmoid activation controlling when drop $n$ triggers.
\end{itemize}

\textbf{Physics constraints} (from the Laws of Thermodynamics):
\begin{itemize}[noitemsep]
  \item \textbf{LoT1 (energy conservation):} predicted foundation frequency must match
        analytical frequency --- enforced via $\mathcal{L}_{\text{freq}}$.
  \item \textbf{LoT2 (non-negative dissipation):} stiffness can only decrease, never
        ``gain'' --- enforced via $\mathcal{L}_{\text{1mon}}$, $\mathcal{L}_{\text{2mon}}$,
        $\mathcal{L}_{\text{pos}}$.
\end{itemize}

\textbf{Sources:}
\begin{itemize}[noitemsep]
  \item Gorini, D.N. (2025). ``A unified thermodynamic-based macroelement for
        geotechnical systems.'' \textit{Computers and Geotechnics}, 179, 106965.
        --- \href{https://www.sciencedirect.com/science/article/pii/S0266352X2200547X}{[Link]}
  \item PostTyphoon PDF, Section 4.2.1, Figures 14--16, pp.\,30--33.
\end{itemize}

% ------------------------------------------------------------
\subsection{Step 2: Prepare the Training Data}
% ------------------------------------------------------------

Two data sources feed the model:

\begin{enumerate}[label=\textbf{\alph*)}]
  \item \textbf{High-fidelity FEM simulations ($\sim$300 cases):}
  \begin{itemize}[noitemsep]
    \item OpenSees 3D FEM with MCP bounding-surface plasticity model for clay.
    \item Dynamic pile pushover analyses across the parameter ranges in Table~\ref{tab:params}.
    \item Extract load--displacement ($p$--$y$) and moment--rotation ($m$--$\theta$) curves
          $\Rightarrow$ determine $H^0_{ij}$ (initial tangent) and $H^n_{ij}$ (secant drops).
    \item Also extract participating mass matrices $m_{ij}$ from strain fields.
  \end{itemize}

  \item \textbf{Scaled physical experiments ($\sim$30 cases):}
  \begin{itemize}[noitemsep]
    \item 1:50 scale monopile-turbine models in Shenzhen marine clay.
    \item Impact hammer tests at each typhoon loading stage to measure $f_0^{\text{OWT}}$.
    \item Provides realistic load--$f_0$ data pairs for final model training.
  \end{itemize}
\end{enumerate}

\begin{table}[h]
\centering
\caption{Parameter ranges for data generation (from PostTyphoon PDF, Table 2, p.\,25).}
\label{tab:params}
\begin{tabular}{ll}
\toprule
\textbf{Parameter} & \textbf{Range} \\
\midrule
Turbine size & 5--16 MW \\
Pile diameter $D$ & 4--10 m \\
Pile length $L$ & 25--60 m \\
Pile bending stiffness $EI$ & 30--150 GNm$^2$ \\
Clay strength $s_u$ & 20--130 kPa \\
Clay stiffness $G_{\max}$ & 10--50 MPa \\
Typhoon return periods & 5, 10, 50-year \\
Water depth & 15--60 m \\
\bottomrule
\end{tabular}
\end{table}

\textbf{Sources:}
\begin{itemize}[noitemsep]
  \item PostTyphoon PDF, Sections 2.1.1--2.1.2 and 4.1.1--4.1.2, pp.\,19--29.
  \item Kato et al. (2023), Ocean Eng., 114338 --- prior 108 FEM simulations.
\end{itemize}

% ------------------------------------------------------------
\subsection{Step 3: Design the Self-Calibration Module (Slot-Attention)}
% ------------------------------------------------------------

This module learns to \textit{auto-calibrate} the TIM stiffness matrices from basic
soil/pile parameters, eliminating case-by-case manual calibration.

\textbf{Architecture} (based on Locatello et al., 2020, Slot Attention):
\begin{enumerate}[noitemsep]
  \item \textbf{21 learnable slots:}
    \begin{itemize}[noitemsep]
      \item Slot 1 $\rightarrow$ generates $H^0_{ij}$ (initial 3DOF stiffness matrix).
      \item Slots 2--21 $\rightarrow$ each generates one $H^n_{ij}$ drop matrix + its
            sigmoid activation $w_n(t)$.
    \end{itemize}
  \item \textbf{Cross-attention layer:} maps encoder tokens (soil/pile properties)
        to each slot.
  \item \textbf{Self-attention layer:} shares context between slots to ensure each
        specializes in a distinct stiffness drop (no overlap).
  \item \textbf{MLP decoders:} per-slot shallow networks that output the actual
        matrix values $H^0_{ij}$, $H^n_{ij}$, $w_n(t)$.
\end{enumerate}

\textbf{Pre-training loss:}
\begin{equation}
  \mathcal{L}_{\text{total}}^{\text{pre}} = \alpha\,\mathcal{L}_{\text{data}}^{\text{pre}}
  + \beta\,\mathcal{L}_{\text{freq}}
  + \gamma\,(\mathcal{L}_{\text{1mon}} + \mathcal{L}_{\text{2mon}})
  + \delta\,\mathcal{L}_{\text{pos}}
\end{equation}

where:
\begin{itemize}[noitemsep]
  \item $\mathcal{L}_{\text{data}}^{\text{pre}}$: MSE between predicted and FEM
        stiffness matrices.
  \item $\mathcal{L}_{\text{freq}}$: predicted vs.\ analytical foundation frequency
        (LoT1).
  \item $\mathcal{L}_{\text{1mon}}$: monotonic decrease of $H^n_{ij}$ drops (LoT2).
  \item $\mathcal{L}_{\text{2mon}}$: monotonic increase of $w_n(t)$ activations (LoT2).
  \item $\mathcal{L}_{\text{pos}}$: non-negative $H_{ij}(t)$ at all times (LoT2).
\end{itemize}

\textbf{Sources:}
\begin{itemize}[noitemsep]
  \item Locatello et al. (2020). ``Object-centric learning with slot attention.''
        \textit{NeurIPS}, 33, 11525--11538.
  \item PostTyphoon PDF, Section 4.2.1, Figures 15--16, pp.\,30--33.
\end{itemize}

% ------------------------------------------------------------
\subsection{Step 4: Build the Full Transformer Architecture}
% ------------------------------------------------------------

The full model follows a classic \textbf{encoder--decoder transformer}:

\begin{enumerate}[noitemsep]
  \item \textbf{Encoder:}
    \begin{itemize}[noitemsep]
      \item Tokenize fixed inputs (soil: $G_{\max}$, $s_u$; pile: $D$, $L$, $EI$)
            and load time-series with sinusoidal positional encodings.
      \item Self-attention layers process the embeddings.
      \item Feed into the pre-trained self-calibration module $\Rightarrow$ outputs
            updated $H_{ij}(t)$ at each load step.
      \item An additional self-attention layer encodes the stiffness information.
    \end{itemize}

  \item \textbf{Decoder:}
    \begin{itemize}[noitemsep]
      \item Masked self-attention (prevents future-data leakage).
      \item Cross-attention to encoder outputs: relates $H_{ij}(t)$ and load info
            to previous $f_0^{\text{OWT}}(t)$ outputs.
      \item Feed-forward network + dense linear layer $\Rightarrow$ predicts both
            $H_{ij}(t)$ and $f_0^{\text{OWT}}(t)$ sequences.
    \end{itemize}
\end{enumerate}

\textbf{Training loss:}
\begin{equation}
  \mathcal{L}_{\text{total}} = \alpha\,\mathcal{L}_{\text{data}}
  + \beta\,\mathcal{L}_{\text{freq}}
  + \gamma\,(\mathcal{L}_{\text{1mon}} + \mathcal{L}_{\text{2mon}})
  + \delta\,\mathcal{L}_{\text{pos}}
\end{equation}

where $\mathcal{L}_{\text{data}}$ now uses NRMSE + MAPE on $f_0^{\text{OWT}}$.
Experimental data receive higher loss weight than analytical data.

\textbf{Source:} PostTyphoon PDF, Section 4.2.2, Figure 17, pp.\,33--35.

% ------------------------------------------------------------
\subsection{Step 5: Attention Roll-Out for XAI (Opening the Black Box)}
% ------------------------------------------------------------

This is the key innovation for Objective 1 --- \textbf{quantifying the correlation chain}.

\begin{enumerate}[noitemsep]
  \item \textbf{Record attention weights} from every multi-head attention block during
        inference (encoder, slot-attention, decoder).
  \item \textbf{Mean-head aggregation:} average attention weights across all heads
        at each layer to get smooth $Q \times K$ matrices.
  \item \textbf{Attention heat maps:}
    \begin{itemize}[noitemsep]
      \item Rows = encoder keys (soil params, load steps).
      \item Columns = decoder queries (slot activations, $f_0$ time steps).
      \item Visualize as heat maps with diverging colormaps.
      \item Generate separate maps for the slot-attention module and the full transformer.
    \end{itemize}
  \item \textbf{Attention roll-out:}
    \begin{itemize}[noitemsep]
      \item Start from the final decoder cross-attention weights.
      \item Propagate backward by multiplying encoder self-attention matrices across
            all layers.
      \item Result: an \textit{attribution vector} showing end-to-end contribution of
            each input token to the final $f_0^{\text{OWT}}$ prediction.
      \item Visualize as a bar chart aligned with load inputs and fixed parameters.
    \end{itemize}
  \item \textbf{Physics validation:} Compare attention patterns against known physical
        expectations (e.g., peak loads should get high attention, $K_R$ should dominate
        $f_0$ sensitivity per Bhattacharya 2019).
\end{enumerate}

\textbf{Sources:}
\begin{itemize}[noitemsep]
  \item PostTyphoon PDF, Sections 2.3.2 and 4.3.2, pp.\,23--24 and 36--37.
  \item Bhattacharya, S. (2019). \textit{Design of foundations for offshore wind
        turbines}. John Wiley \& Sons. --- showed $K_R$ dominates $f_0$ shifts.
\end{itemize}

% ------------------------------------------------------------
\subsection{Step 6: Testing, Validation, and Iteration}
% ------------------------------------------------------------

\begin{enumerate}[noitemsep]
  \item \textbf{Train/test split:} 85\%/15\% random split for both FEM and experimental
        data. 3 out-of-domain experiments reserved for final validation.
  \item \textbf{Self-calibration module testing:} MSE and $R^2$ on unseen FEM stiffness
        matrices; pass if MAPE $< 10\%$ and bias $< 5\%$.
  \item \textbf{Full model testing:} MSE and $R^2$ on unseen $f_0^{\text{OWT}}$ data;
        target MAPE $< 1\%$ (since even 10\% $f_0$ shift is critical).
  \item \textbf{Validation data:}
    \begin{itemize}[noitemsep]
      \item Centrifuge experiments from Lai et al.\ (2020) for foundation stiffness.
      \item Field measurements from a 3.6 MW deployed OWT during a real storm.
      \item 3 out-of-domain scaled experiments beyond Table~\ref{tab:params} ranges.
    \end{itemize}
  \item \textbf{Uncertainty quantification:} Add Gaussian noise to inputs; if CoV $> 10\%$,
        refine attention regularization or physics losses.
  \item \textbf{Iterate:} Use attention heat maps to identify weak points, prune
        low-attention inputs, increase sampling in high-attention load windows.
\end{enumerate}

\textbf{Source:} PostTyphoon PDF, Sections 2.3.1 and 4.3.1, pp.\,22--23 and 35--36.

% ============================================================
\section{Summary: The Big Picture}
% ============================================================

\begin{center}
\begin{tabular}{cp{10cm}}
\toprule
\textbf{Step} & \textbf{What You Do} \\
\midrule
1 & Learn TIM physics: how $H_{ij}(t)$ degrades via yield surfaces \\
2 & Generate training data: $\sim$300 FEM + $\sim$30 experiments \\
3 & Pre-train slot-attention module: learns $H^0_{ij}$, $H^n_{ij}$, $w_n(t)$ from soil/pile params \\
4 & Train full transformer: encoder (loads + self-calibration) $\to$ decoder ($f_0^{\text{OWT}}$) \\
5 & Attention roll-out: quantify each link in the soil $\to$ foundation $\to$ $f_0$ chain \\
6 & Test, validate, iterate using XAI insights \\
\bottomrule
\end{tabular}
\end{center}

% ============================================================
\section{Key References for Objective 1}
% ============================================================

\begin{enumerate}[noitemsep]
  \item Cheng et al.\ (2023). Dynamic response analysis of monopile OWTs considering
        stiffness degradation of clay. \textit{Comput.\ Geotech.} ---
        \href{https://www.sciencedirect.com/science/article/pii/S0266352X2200547X}{Link}

  \item Kato, Bhattacharya \& Wang (2023). Evaluation of post-storm soil stiffness
        degradation effects on OWTs in clay. \textit{Ocean Eng.}, 114338 ---
        \href{https://www.sciencedirect.com/science/article/pii/S0029801823007229}{Link}

  \item Ding, Chian et al.\ (2023). Wind-wave combined effect on soil-monopile-OWT.
        \textit{Comput.\ Geotech.} ---
        \href{https://www.sciencedirect.com/science/article/pii/S0266352X2200461X}{Link}

  \item Xu, Huang et al.\ (2025). OWT monopiles under real wind and wave:
        cyclic softening of clays. \textit{Ocean Eng.} ---
        \href{https://www.sciencedirect.com/science/article/pii/S0029801825001647}{Link}

  \item Cheng, Xing et al.\ (2025). Rotation accumulation and natural frequency
        degradation for OWT in clays. \textit{Ocean Eng.} ---
        \href{https://www.sciencedirect.com/science/article/pii/S0029801824030993}{Link}

  \item Gorini, D.N.\ (2025). A unified thermodynamic-based macroelement. \textit{Comput.\
        Geotech.}, 179, 106965.

  \item Locatello et al.\ (2020). Object-centric learning with slot attention.
        \textit{NeurIPS}, 33, 11525--11538.

  \item Bhattacharya, S.\ (2019). \textit{Design of foundations for offshore wind
        turbines}. Wiley.

  \item Shadlou \& Bhattacharya (2016). Dynamic stiffness of monopiles. \textit{Soil
        Dyn.\ Earthq.\ Eng.}, 88, 15--32.

  \item Lai et al.\ (2020). Centrifuge modeling of cyclic lateral behavior of
        monopiles in soft clay. \textit{Ocean Eng.}, 200, 107048.
\end{enumerate}

% ============================================================
\section{Companion Code}
% ============================================================

A simplified but comprehensive Python demo is provided in the companion file:

\begin{center}
  \texttt{objective1\_XAI\_transformer\_demo.py}
\end{center}

This demo implements a minimal version of the full architecture (slot-attention
self-calibration module + encoder-decoder transformer + physics losses +
attention roll-out) using PyTorch, with synthetic data mimicking the real
parameter ranges. It is meant as a \textbf{proof-of-concept} to understand each
component before scaling to real FEM/experimental data.

\end{document}
